\chapter{كم مدخلًا نحتاج}
% full version: chapters/19_dendrites.tex

للعصبون زوائد مدخلة متفرعة: التغصنات. عند بعض العصبونات هي شجرة كاملة بآلاف نقاط التماس، وعند أخرى فرع أو فرعان أو لا شيء أصلًا. بالنسبة إلى المهندس، عدد المداخل معامل في الدائرة، وهو مختار بحسب المهمة.

\section*{من الصفر إلى مئة ألف}

مستشعرات اللمس والألم بلا تغصنات إطلاقًا: السلك يمتد ببساطة من الجلد إلى الحبل الشوكي، والمستشعر قائم في طرفه نفسه. إنه خط بمستشعر في طرفه البعيد.

\pencilfig{19}{\pencilart{19}}{مداخل كثيرة تعني جامعًا؛ ومدخل واحد ضخم يعني مرحّلًا سريعًا ودقيقًا.}

عصبونات الشم، والعصبونات التي تنقل الإشارة من مستشعرات العين والأذن، لها تغصن واحد. مدخل واحد ومخرج واحد: مكرِّر، أو عازل (\foreignlanguage{english}{buffer}).

وكواشف اتجاه الصوت لها تغصنان يتجهان في اتجاهين مختلفين، واحد لكل أذن. بتوزيع المدخلين على فرعين مختلفين يجعل العصبون ”و“ أكثر صرامة: فحصتا تيار في فرع واحد كانتا ستعيق إحداهما الأخرى.

والعصبونات الصغيرة في دائرة تعلم الحركات لها نحو أربعة تغصنات قصيرة، كل منها يلتقط مدخله. هذه العصبونات بعشرات المليارات، وكل منها لا يجمع إلا أربع إشارات. لكنها معًا تفرد المدخل في ملايين التراكيب، ويختار منها عصبون الخرج الكبير ذو المئة ألف مدخل ما يلزمه.

وأخيرًا، العصبونات ذات الأشجار الضخمة وآلاف المداخل: جامعات ومصنِّفات، تعمل فروعها المنفردة تقريبًا كدوائر صغيرة مستقلة.

\section*{لماذا هكذا}

الشحن هو ما يحسم كل شيء. العصبون الصغير ذو المداخل القليلة يُشحن من وصلة كبيرة واحدة في أجزاء من جزء من ألف من الثانية: عصبون كهذا دقيق في الزمن وينقل الإشارة موثوقًا. أما الشجرة الكبيرة فلن تشحنها وصلة واحدة أبدًا: ستتسرب الشحنة قبل ذلك. تحتاج إلى عشرات الإشارات الواصلة معًا، لكنها في المقابل تستطيع جمعها ومقارنتها. مداخل قليلة ووصلات كبيرة: للتوقيت الدقيق. مداخل كثيرة: للعدّ.

\pencilfig{128}{%
% charging: a small neuron reaches threshold from one large synapse at once; a big tree barely moves
% from one input and needs many inputs arriving together
\foreach \dx/\t in {0/{مداخل قليلة}, 4.3/{شجرة كبيرة}}{
  \draw[pen] (\dx,0) -- (\dx+3.6,0);
  \draw[pcGray, densely dashed, line width=0.5pt] (\dx,0.9) -- (\dx+3.6,0.9);
  \node[font=\small\itshape, text=pcGray, anchor=south] at (\dx+1.8,1.75) {\t};}
\node[lbl, anchor=east] at (-0.05,0.9) {العتبة};
% small neuron
\draw[pcBlue, line width=2pt] (0.6,-0.15) -- (0.6,-0.6);
\draw[penlite=pcBlue] (0,0) -- (0.6,0) -- plot[domain=0.6:0.8, samples=12] (\x,{1.3*(1-exp(-(\x-0.6)/0.18))}) -- (0.8,1.6) -- (0.84,0) -- (3.5,0);
\node[lbl, text=pcRed, anchor=west] at (0.85,1.45) {نقرة فورًا};
\node[lbl, text=pcBlue, anchor=north] at (0.6,-0.62) {مدخل كبير واحد};
% big tree
\draw[pcBlue, line width=0.6pt] (4.8,-0.15) -- (4.8,-0.45);
\foreach \s in {6.15,6.2,6.25,6.3,6.35,6.4,6.45,6.5}{\draw[pcBlue, line width=0.6pt] (\s,-0.15) -- (\s,-0.45);}
\draw[penlite=pcBlue] (4.3,0) -- (4.8,0) -- plot[domain=4.8:6.15, samples=20] (\x,{0.16*((\x-4.8)/0.15)*exp(1-(\x-4.8)/0.15)}) -- plot[domain=6.15:6.62, samples=14] (\x,{1.25*(1-exp(-(\x-6.15)/0.4))}) -- (6.62,1.6) -- (6.66,0) -- (7.9,0);
\node[lbl, anchor=south] at (4.95,0.2) {لا شيء تقريبًا};
\node[lbl, text=pcBlue, anchor=north] at (6.33,-0.47) {كثيرة معًا};
}{العصبون الصغير يُشحن من وصلة كبيرة واحدة في أجزاء من جزء من ألف من الثانية وينقر فورًا. أما الشجرة الكبيرة فلا تكاد تتغير من مدخل واحد: تحتاج إلى عشرات الإشارات الواصلة معًا.}

وهكذا حصلنا على فرز ثالث للعصبونات: بحسب المادة، وبحسب الجهاز، وبحسب عدد المداخل. وكل فرز يرى الآلة نفسها من جانب جديد.

\fullversion{19}
